\documentclass{subfiles}
\begin{document}
    Let $X$ be the domain of our problem, and let $x \in X$ be a solution.
        It is very useful to define the so called \emph{neighborhood} $N(x)$ of $x$,
        which represents the set of points in the proximity of $x$.
        The idea behind the notion of neighborhood is simple:
        since $x$ is a solution, all the surrounding points are solutions as well,
        and there could be one that is better.

    From the notion of neighborhood, we can derive the notion of local optimum as follows:
        if $(X, c)$ is an instance of a problem, where $X$ is the domain of the problem,
        and $c$ the const functions to minimize (resp. maximize), 
        and $N(x)$ is the neighborhood of a solution $x$,
        we say that $x$ is a local optimum if
        $$
            \forall x' \in N(x), c(x) \le c(x').
        $$
\end{document}
